% !PRINTED 272
\weberpage{272}
Wir nehmen an, dass wir auf irgend eine Weise eine Reihe von $m + 1$ stetigen Functionen herstellen können
\weq{1}{f(x), f_1(x), f_2(x) \ldots f_m(x), \quad (\mathfrak{R})}
denen folgende Eigenschaften zukommen:

1. Von den Functionen $f_\nu$ sollen im Intervall $\Delta$ nicht zwei auf einander folgende zugleich verschwinden.

2. Die letzte von ihnen, $f_m(x)$, soll im Intervall $\Delta$ überhaupt nicht verschwinden, also ein unveränderliches Zeichen behalten.

3. Wenn ein mittleres Glied, etwa $f_\nu(x)$ für irgend ein $x$ im Intervall $\Delta$ verschwindet, so sollen für dieses $x$ die beiden angrenzenden Functionen $f_{\nu-1}(x)$ und $f_{\nu+1}(x)$ entgegengesetztes Vorzeichen haben.

4. Wenn $f(x)$ im Intervall $\Delta$ verschwindet, so soll $f_1(x)$ für diesen Werth von $x$ dasselbe Vorzeichen haben wie $f'(x)$.

Eine solche Functionenreihe ($\mathfrak{R}$) wollen wir eine \gesperrt{Sturm'sche Kette} nennen. Wie man sie bilden kann, werden wir später sehen; zunächst sollen aus der Definition Folgerungen gezogen werden.

Für jeden Werth von $x$, für den keine der Functionen von ($\mathfrak{R}$) verschwindet, hat jede dieser Functionen ein bestimmtes Vorzeichen. Wir zählen einen Zeichenwechsel, so oft beim Durchlaufen der Kette von links nach rechts auf ein positives Glied ein negatives oder auf ein negatives Glied ein positives folgt. Die Anzahl der so gezählten Zeichenwechsel (Variationen) für einen bestimmten Werth von $x$ wollen wir mit $V(x)$ bezeichnen. Wenn ein mittleres Glied $f_\nu(x)$ verschwindet, so haben wir nach 3. beim Uebergang von $f_{\nu-1}(x)$ zu $f_{\nu+1}(x)$ einen Zeichenwechsel zu zählen.

Sind $\alpha$, $\beta$ zwei Werthe, für die keine der Functionen $f_\nu$ verschwindet, und lassen wir nun $x$ stetig wachsen von $\alpha$ bis $\beta$, so wird eine Aenderung in der Zahl der Zeichenwechsel nur dann eintreten können, wenn eine der Functionen von ($\mathfrak{R}$) ihr Zeichen wechselt, also durch den Werth Null hindurchgeht (wegen der vorausgesetzten Stetigkeit).

Ist dies aber eine mittlere Function, so ändert sich die Zahl der Zeichenwechsel nicht (nach 3). Denn in
\[
f_{\nu-1}, \; f_\nu, \; f_{\nu+1}
\]
%CONT>>
